\documentclass[10pt]{article}
\usepackage[margin=20mm]{geometry}
\usepackage{amsmath,amssymb,mathtools,fontspec,unicode-math,xurl,hyperref}
\IfFontExistsTF{DejaVu Serif}{\setmainfont{DejaVu Serif}}{\setmainfont{Latin Modern Roman}}
\setmathfont{latinmodern-math.otf}
\setlength{\parindent}{0pt}\setlength{\parskip}{6pt}\sloppy
\title{Technical review of Quasi-coincidence of cluster structures on positroid varieties}
\author{AI-assisted technical review}
\date{}
\begin{document}
\maketitle
Abstract: The bound accepted-author manuscript is examined through the disconnected reduction, the categorical comparison, and the twist argument. A local inconsistency between ordinary generic-fibre dimension and normalized module rank is confirmed by an explicit circle-algebra module. It is assessed as Minor: the subsequent constructions use the intended vertex-normalized modules, and no damage to the main quasi-coincidence argument is established. Independent determinant and sign checks support the inspected examples. Actual inspection of the source figures and report images remains blocked.\par
\section*{Source, definitions, and scope}
The reviewed work is Matthew Pressland, Quasi-coincidence of cluster structures on positroid varieties, institutional accepted manuscript in Glasgow deposit 368020, deposited 5 December 2025 [paper]. The journal identity is DOI 10.4171/jems/1736 (https://doi.org/10.4171/jems/1736). The exact PDF contains the 59-page article followed by the repository notice on PDF page 60. Article page numbers below are also PDF page numbers. No identity with the publisher PDF or an arXiv edition is assumed.\par
All 60 pages of extracted text, including references and the repository notice, were actually read. This is not a claim to have independently proved all cited results, or to have inspected the ten figures as images. The principal reconstructed arguments are the Segre-product reduction in Section 4, the index computation and derived comparison in Section 6, and the categorical twist deduction in Section 7.\par
The ring is the coordinate ring of the cone over an open positroid variety, with the necklace Plücker coordinates inverted. Definition 3.4 requires the frozen-monomial condition, an isomorphism after coefficient specialization, and compatibility of the \(\widehat y\) variables. Equality of the sets of ordinary cluster variables is a stronger assertion and is not the theorem. Connectedness is imposed before using the boundary algebra and its categorifications; the disconnected case is treated geometrically first. Remark 3.8 also explicitly excludes the stronger conjecture concerning every relabelled cluster structure.\par
\section*{P1: inconsistent normalization of module rank}
\subsection*{Location and decisive calculation}
Definition 5.5, page 31, sets \(Z=\mathbb C[[t]]\), \(K=\mathbb C((t))\), and displays \[\operatorname{rank}(M)=\dim_K(K\otimes_ZM).\] It then states, for the dimer algebra \(A_D\), the normalized formula \[\operatorname{rank}(M)=\frac{1}{\#Q_0}\operatorname{rank}_Z(M).\] The first displayed expression is instead exactly \(\operatorname{rank}_Z(M)\) for a finite free \(Z\)-module. On the same page, the modules \(M_I\) of Definition 5.7, having one copy of \(Z\) at every circle-quiver vertex, are asserted to have rank one. These statements cannot all use the first displayed definition. Both layout and raw extraction of the exact bound page give this normalization mismatch.\par
An explicit admissible example removes any reliance on an informal dimension count. Take \(C_{2,4}\) and \(I=\{1,2\}\). On the total module \(Z^4\), the sums of the clockwise and anticlockwise arrow maps specified in Definition 5.7 are \[X=\begin{pmatrix}
0&0&0&1\\ t&0&0&0\\0&t&0&0\\0&0&1&0
\end{pmatrix},
\qquad
Y=\begin{pmatrix}
0&1&0&0\\0&0&1&0\\0&0&0&t\\t&0&0&0
\end{pmatrix}.\] They satisfy \[XY=YX=tI_4,\qquad Y^2=X^2.\] Thus they satisfy the circle-algebra relations \(xy=yx\) and \(y^k=x^{n-k}\) with \(k=2,n=4\). Each vertex space is free of rank one over \(Z\), so this is the stipulated Cohen–Macaulay module \(M_{\{1,2\}}\). Its generic fibre is \(K^4\), of \(K\)-dimension four. Its normalized vertex rank is one. The two notions differ on an object satisfying all the module assumptions, without invoking any unproved classification of cluster variables.\par
These matrix relations and dimensions were checked by a reviewer-written symbolic script using the existing Python/SymPy runtime, with final exit code 0. The displayed algebra also provides a direct proof.\par
\subsection*{Impact, correction, and limits}
The assessment is Minor. The defect concerns the definition and its compatibility with the subsequently stated rank-one objects. The inspected constructions explicitly specify their vertex modules, perfect matchings, and cluster-tilting summands; they do not infer a new homological vanishing statement from the erroneous unnormalized dimension. The subsequent paragraph already states the intended normalization for \(A_D\). No counterexample to Theorems 6.16 or 7.2 is established.\par
A precise local correction is to distinguish ordinary \(Z\)-rank from the normalized rank of the orders used in the paper. When an order has generic fibre \(\operatorname{Mat}_m(K)\), the normalized rank is \[\frac{1}{m}\dim_K(K\otimes_ZM),\] equivalently the multiplicity of its generic simple module. For \(C_{k,n}\) this uses \(m=n\); for \(A_D\) it uses the displayed factor \(\#Q_0\). A single such normalization should not be asserted for an arbitrary finite free \(Z\)-algebra without specifying its generic fibre. The correction restores the rank of the explicit module to one. Author adoption has not been requested or observed.\par
\section*{Checks of the principal argument}
\subsection*{Disconnected reduction and frozen units}
The cut through a boundary face produces two quivers meeting at a single frozen vertex, with no arrows between their mutable parts. Mutations in the two parts therefore commute. The coordinate-ring comparison is the graded Segre product, not the unrestricted tensor product. The sign in Corollary 4.3 is \[(-1)^{k_2|I\cap S_1^+|}.\] Directly, the only inversions when concatenating the ordered subsets \(I\subset S_1\) and \(J\subset S_2\) are pairs with the first member in \(S_1^+\). There are \(|J|\,|I\cap S_1^+|=k_2|I\cap S_1^+|\) such pairs. An exact enumeration checked this equality in 3356 admissible finite cases; the inversion argument proves the general formula.\par
The degree identities in Propositions 4.13 and 4.16 ensure that both terms of an exchange relation carry the same sign and auxiliary frozen power. The correction factor in Theorem 4.17, \(p(q/x_*)^{\deg x}\), is permitted because frozen variables are inverted. Moreover, each exchange column has zero pairing with every cluster grading. Thus the \(\widehat y\) variables have zero standard degree and zero relevant weight degree: the auxiliary frozen powers and signs disappear in their comparison. This independently explains why the third condition of Definition 3.4 survives the disconnected reduction. It does not extend the result to coordinate rings with noninvertible frozen variables.\par
\subsection*{Stable projection, shifts, and the categorical criterion}
The operative comparison is \[(\operatorname{ginj}\operatorname{CM}(B),T^{\rm tgt})
\longrightarrow
(\operatorname{gproj}\operatorname{CM}(B),T^{\rm src}).\] These are the categories and objects constructed in Theorems 5.12–5.15 and in diagrams (5.8), (6.4), and (7.1). A reversed pairing of superscripts in the closing paragraph of the proof of Theorem 6.16 should be read against these explicit constructions; it does not supply a different categorical proof.\par
The crucial equality of matching classes in Lemma 6.10 is made only after projection to the Grothendieck group of the stable endomorphism algebra. Remark 6.11 explicitly rejects the corresponding equality before projection. Using the projected equality, Corollary 6.8 gives \[\operatorname{ind}\!\left(\Omega(e_jAe)^\vee\right)=-[Ae_j].\] The relevant object is rigid in the singularity category, and the cited uniqueness of rigid objects with a given index yields \[\Omega(e_jAe)^\vee\simeq\Sigma(eAe_j).\] Since \(\Omega=\Sigma^{-1}\) there, this gives \(\Sigma^2(eAe_j)\simeq(e_jAe)^\vee\), with the shift direction as in Theorem 6.4. The suspension closure of the mutation class in Proposition 6.3 then supplies the reachability required by Fraser–Keller’s criterion.\par
Proposition 6.12 additionally supplies the lift to bounded homotopy categories. Its argument identifies the transformed complex with \(F'T^{\rm tgt}=A^\vee\) and uses the finite global dimension and tilting properties. The cluster-character comparison on the initial summands is proved separately in Proposition 6.14, using the twist formula and reachability. Thus a bare derived equivalence is not being treated as sufficient for quasi-coincidence. The foundational categorification results, the green-to-red sequence theorem, rigid-index uniqueness, and Fraser–Keller’s criterion remain cited inputs rather than independently reproved results.\par
The twist proof then replaces the comparison by its inverse shift. Equation (7.2) contributes the projective denominator, and \(\Sigma^{-1}T^{\rm tgt}\simeq\Sigma T^{\rm src}\) supplies the required mutation class. This checks the categorical route to Theorem 7.2 without assuming in advance that the two original variable sets coincide.\par
\subsection*{Independent check on the stated open positroid}
For Example 2.18 and formula (6.9), take the row space of \[A=\begin{pmatrix}
2&1&1&0&0&0&0\\
3&1&2&1&0&1&1\\
3&0&0&0&1&1&2
\end{pmatrix}.\] The matrix has rank three; its five prescribed minors \(\Delta_{234},\Delta_{456},\Delta_{457},\Delta_{467},\Delta_{567}\) vanish. Every frozen minor in both necklaces is nonzero, so this is a point of the stated open positroid, not merely an arbitrary matrix. Exact calculation gives \[\Delta_{157}=-2,\quad \Delta_{357}=-1,\quad
\Delta_{167}=2,\quad\Delta_{367}=1,\] and hence \(\Delta_{157}=\Delta_{357}\Delta_{167}/\Delta_{367}\) as claimed.\par
A preliminary ambient determinant test failed because of the sign of the term involving \(\Delta_{567}\). With ordered columns, the universal relation is \[\Delta_{357}\Delta_{167}-\Delta_{157}\Delta_{367}
=-\Delta_{137}\Delta_{567}.\] The manuscript uses the other sign before setting this term to zero. On its specified positroid \(\Delta_{567}=0\), so this does not refute the displayed on-variety equality or formula (6.9). The failed script, exception, and receipt are retained, alongside the corrected test that explicitly checks the on-variety hypotheses. No unrestricted matrix is presented as a counterexample to the positroid result.\par
\section*{Outcome and limitations}
The supported finding is the local rank-normalization inconsistency P1, assessed as Minor. No core theorem failure has been established. Full text reading does not certify all external categorical theorems, the exact combinatorics drawn in the figures, or the whole mutation class through computation.\par
The repository, publisher, author publication list, and bounded searches did not identify a separate supplement or executable repository for this exact paper. Corrections listed for other works on the author’s page are not assigned to this paper. No author code was executed.\par
Item | State | Scope\par
P1 | Minor | Rank definition versus explicit vertex modules\par
Independent checks | Passed, final run | Circle-algebra relations, admissible minors, Segre signs\par
Repair/adoption | Local formula checked | No author adoption observed\par
Scientific decision | Pending | No independent or human acceptance\par
\section*{Supporting checks}
> 00000TOKEN.\par
Paper ID: `M25-0357`.\par
Scientific reasoning and writing are by this assistant session; native tools perform mechanical calculations, compilation, parsing, and rendering. Controller UI observations are attributed to the controller, with backend equivalence unverified. Authorized image reads return `EMFILE`. Rendering is not viewing, and no actual image inspection is claimed. Delivery completion remains withheld.\par
9 M. Pressland, Quasi-coincidence of cluster structures on positroid varieties, accepted-author manuscript, Glasgow deposit 368020, https://eprints.gla.ac.uk/368020/, deposited 5 December 2025. EMS Press, DOI 10.4171/jems/1736, https://ems.press/journals/jems/articles/14299327, consulted 27 October 2026. M. Pressland, publication list, https://mdpressland.github.io/papers.html, consulted 27 October 2026.\par
\end{document}
